\begin{solapartat}
\punts{0,1} El potencial s'expressa com a:
\[ V = k\,\frac{q}{r} \;\Rightarrow\; r = k\,\frac{q}{V} \]

\punts{0,4} $r_1 = k\,\dfrac{Q}{V_1} = 9\times 10^{9}\,\dfrac{1\times 10^{-9}}{3} = 3\ \mathrm{m}$, $r_2 = k\,\dfrac{Q}{V_2} = 9\times 10^{9}\,\dfrac{1\times 10^{-9}}{6} = 1,5\ \mathrm{m}$,
\[ r_3 = k\,\frac{Q}{V_3} = 9\times 10^{9}\,\frac{1\times 10^{-9}}{9} = 1\ \mathrm{m} \text{ i } r_4 = k\,\frac{Q}{V_4} = 9\times 10^{9}\,\frac{1\times 10^{-9}}{12} = 0,75\ \mathrm{m}. \]

\punts{0,5} Per a la representació correcta de la projecció de les superfícies equipotencials, si aquestes no són circumferències, es restaran 0,25 punts. Si no estan a la distància correcta, es restaran 0,25 punts.

\figura[0.5]{sol_fig1.png}

\punts{0,25} La distància entre les superfícies equipotencials no és constant. Són, respectivament, 1,5~m, 0,5~m i 0,25~m.
\end{solapartat}

\begin{solapartat}
% DUBTE: la pauta oficial (pau_fisi22sp.pdf, pàg. 5) no conté cap text ni puntuació
% específics per a aquest apartat b); només tracta l'apartat a). La figura de la
% pauta (sol_fig1.png) sí que inclou les 8 línies de camp dibuixades radialment des
% de la càrrega, tallant perpendicularment les circumferències equipotencials, però
% no hi ha cap enunciat del criteri de correcció ni la resposta explícita sobre
% l'angle (que, per tractar-se de línies de camp i equipotencials, és de 90 graus).
\textit{La pauta oficial no inclou cap text ni puntuació per a aquest apartat.}
\end{solapartat}
